\begingroup
\setlength{\tabcolsep}{4.5pt}
\renewcommand{\arraystretch}{1.28}
\begin{longtable}{@{}>{\raggedright\arraybackslash}p{.15\textwidth}>{\raggedright\arraybackslash}p{.31\textwidth}>{\raggedright\arraybackslash}p{.25\textwidth}>{\raggedright\arraybackslash}p{.23\textwidth}@{}}
\caption{Secteurs topologiques et virtuels : invariants natifs de fusion, de tressage et de persistance.}\label{tab:masas-invariantes-topologicos-virtuales-rev3}\\
\toprule
\textbf{Secteur} & \textbf{Objet et relations} & \textbf{Invariant natif} & \textbf{Portée physique} \\
\midrule
\endfirsthead
\toprule
\textbf{Secteur} & \textbf{Objet et relations} & \textbf{Invariant natif} & \textbf{Portée physique} \\
\midrule
\endhead
\midrule
\multicolumn{4}{r}{\footnotesize Suite à la page suivante.}\\
\endfoot
\bottomrule
\endlastfoot
Fibonacci & $K=\mathbb Z1\oplus\mathbb Z\tau$, $\tau^2=1+\tau$. & $\operatorname{FPdim}(\tau)=\varphi$ dans l’anneau basé. & Les matrices $F/R$, le pentagone, l’hexagone, l’Hamiltonien et le gap relèvent de la réalisation ultérieure. \\
Ising & $K=\mathbb Z\{1,\psi,\sigma\}$, $\psi^2=1$, $\psi\sigma=\sigma$, $\sigma^2=1+\psi$. & $\operatorname{FPdim}(\sigma)=\sqrt{2}$ ; représentation de tressage. & La réalisation de Majorana conserve son propre morphisme physique. \\
$\mathbb Z/6$ & $q_{\epsilon,t}=(-1)^\epsilon\omega_3^t=\zeta_6^k$. & Caractère quadratique et représentation de tressage sur les six secteurs. & L’observation pertinente est topologique, non une masse universelle. \\
Secciones virtuales & Système inverse de prolongements d’horizon $L(s_N)$. & Enregistrement fini, signature, retenue et classe d’obstruction. & L’observable natif est la persistance du prolongement. \\
\end{longtable}
\endgroup
